\section{武周家庭、边地与行政实践} \label{sec:wu-zhou-family-frontier-administration}  \subsection{土地、水与行政} 本段以账本节点 ST-680-LI-XIAN、ST-684-LI-JINGYE、ST-690-ZHOU-FOUNDATION、ST-692-ANXI-RECOVERY、ST-702-BEITING 为时间骨架。土地、河流、山道、草场和城池从不只是政治行动的背景：它们决定户籍如何登记，粮食如何运输，军队在哪些季节可以行动，地方社会又如何在征发与市场之间维持生计。诏令往往以统一的州县、户等和赋役语言描述这些问题，实际执行却受到水路、地形、劳力、疾病与地方关系的限制。制度史可以说明国家试图如何分类，遗址、文书和墓志只能照亮局部实践；两者之间的距离不应被任何整齐的结论填平。  